\documentclass[12pt, a4paper, oneside]{ctexart}
\usepackage{amsmath, amsthm, amssymb, bm, xcolor, tikz, framed, graphicx, hyperref, mathrsfs}

\title{Soild State Physics Homework (Week 5-6)}
\author{Tang Yi \\ [学号已脱敏]}
\usepackage[a4paper, left=2cm, right=2cm, top=2cm, bottom=2cm]{geometry}
\linespread{1.3} % 行距更紧凑
\definecolor{shadecolor}{RGB}{241, 241, 255}

\newcounter{problemname}
\newenvironment{problem}
  {\begin{shaded}\stepcounter{problemname}\par\noindent\textbf{题目 \theproblemname. }}
  {\end{shaded}\par}
\newenvironment{solution}
  {\par\noindent\textbf{题目 \theproblemname 的解答. }}
  {\par}
\newenvironment{note}
  {\par\noindent\textbf{题目 \theproblemname 的注记. }}
  {\par}

\begin{document}
\maketitle
\begin{problem}
证明：在一个周期性势场中，电子的能量是分立的，并且每个能级对应一个布洛赫态。假设 III--V 族化合物中，III 族、V 族原子都是电中性的（$q^* = 0$），求出其电离度 $f_i$。
\end{problem}

\begin{solution}

对于 III--V 化合物，若有效电荷 $q^* = 0$，说明完全共价键合，
即没有电荷转移，因此电离度定义为
\[
f_i = \frac{q^*}{e} = 0.
\]

因此
\[
\boxed{f_i = 0.}
\]
\end{solution}

\begin{problem}
    用林纳德--琼斯势计算 $Ne$ 在体心立方和面心立方结构中的结合能之比值。
\end{problem}

\begin{solution}
    林纳德--琼斯势为
\[
V(r) = 4\epsilon \left[ \left(\frac{\sigma}{r}\right)^{12}
- \left(\frac{\sigma}{r}\right)^6 \right].
\]

晶体总能量可写为
\[
U = 2\epsilon \left( A_{12}\left(\frac{\sigma}{r}\right)^{12}
- A_6 \left(\frac{\sigma}{r}\right)^6 \right),
\]
其中 $A_{12}, A_6$ 为晶格求和常数。

平衡条件：
\[
\frac{dU}{dr} = 0
\Rightarrow r_0 = \sigma \left(\frac{2A_{12}}{A_6}\right)^{1/6}.
\]

代回得最小能量：
\[
U_{\min} = -\epsilon \frac{A_6^2}{2A_{12}}.
\]

已知：
\[
\begin{aligned}
\text{FCC: } & A_6 = 14.45,\quad A_{12} = 12.13 \\
\text{BCC: } & A_6 = 12.25,\quad A_{12} = 10.27
\end{aligned}
\]

则
\[
U_{\text{FCC}} = -\epsilon \frac{(14.45)^2}{2\times 12.13},
\quad
U_{\text{BCC}} = -\epsilon \frac{(12.25)^2}{2\times 10.27}.
\]

比值：
\[
\frac{U_{\text{FCC}}}{U_{\text{BCC}}}
= \frac{(14.45)^2 / 12.13}{(12.25)^2 / 10.27}
\approx 1.09.
\]

因此
\[
\boxed{\frac{U_{\text{FCC}}}{U_{\text{BCC}}} \approx 1.09}
\]
说明 FCC 更稳定。
\end{solution}

\begin{problem}
    对于 $H_2$，从气体的测量得到的林纳德--琼斯势参数为$\epsilon = 50 \times 10^{-13}\,\text{J}$，$\sigma = 2.96\,\text{\AA}$，计算 $H_2$ 结合成面心立方固体分子氢时的结合能（以千焦每摩尔为单位），
    每个氢分子可当做球形来处理。结合能的实验值为 $0.751\,\text{kJ/mol}$，试与计算值进行比较。
\end{problem}

\begin{solution}
    FCC 晶体中 LJ 势总能：
\[
U = 2\epsilon \left( A_{12}\left(\frac{\sigma}{r}\right)^{12}
- A_6\left(\frac{\sigma}{r}\right)^6 \right).
\]

最小值：
\[
U_{\min} = -\epsilon \frac{A_6^2}{2A_{12}}.
\]

代入 FCC 常数：
\[
A_6 = 14.45,\quad A_{12} = 12.13
\]

得
\[
U_{\min} = -\epsilon \frac{(14.45)^2}{2\times 12.13}
\approx -8.6 \epsilon.
\]

代入 $\epsilon = 50\times 10^{-13}\,\text{J}$：
\[
U_{\min} \approx -4.3\times 10^{-12}\,\text{J}.
\]

转为每摩尔：
\[
U = -4.3\times 10^{-12} \times 6.02\times 10^{23}
\approx -2.6\times 10^{12}\,\text{J/mol}
\]

换算为 kJ/mol：
\[
U \approx -2.6\times 10^9\,\text{kJ/mol}
\]

显然远大于实验值 $0.751\,\text{kJ/mol}$。

\[
\boxed{\text{LJ 势严重高估结合能（忽略量子效应与分子结构）}}
\]
\end{solution}

\begin{problem}
    已知一维单原子链，其中第 $j$ 个格波，在第 $n$ 个格点引起的位移 $\mu_{nj}$ 为
\[
\mu_{nj} = a_j \sin(\omega_j t + naq_j + \delta_j)
\]
$\delta_j$ 为任意位相因子。并已知在较高温度下每个格波的平均能量为 $kT$，
试计算每个原子的平均平方位移。
\end{problem}

\begin{solution}
    位移为
\[
\mu_{nj} = a_j \sin(\omega_j t + naq_j + \delta_j)
\]

时间平均：
\[
\langle \sin^2 \rangle = \frac{1}{2}
\Rightarrow \langle \mu_{nj}^2 \rangle = \frac{a_j^2}{2}.
\]

每个模态能量：
\[
E_j = \frac{1}{2}m\omega_j^2 a_j^2 = kT
\Rightarrow a_j^2 = \frac{2kT}{m\omega_j^2}.
\]

因此
\[
\langle \mu_{nj}^2 \rangle = \frac{kT}{m\omega_j^2}.
\]

总平均：
\[
\boxed{
\langle u^2 \rangle = \sum_j \frac{kT}{m\omega_j^2}
}
\]
\end{solution}

\begin{problem}
    讨论 $N$ 个原胞的一维双原子链（相邻原子间距为 $a$），其 $2N$ 个格波解，
    当 $M = m$ 时与一维单原子链的结果一一对应。
\end{problem}

\begin{solution}
    双原子链色散关系：
\[
\omega^2 = \frac{C}{m} \left[ 1 \pm \sqrt{1 - \sin^2\frac{ka}{2}} \right]
\]

当 $M = m$ 时，两支分别为：
\[
\omega = \sqrt{\frac{2C}{m}} \left|\sin\frac{ka}{2}\right|,
\quad
\omega = \sqrt{\frac{2C}{m}} \left|\cos\frac{ka}{2}\right|.
\]

这两支正好拼成单原子链：
\[
\boxed{
\omega = 2\sqrt{\frac{C}{m}}\left|\sin\frac{ka}{2}\right|
}
\]

即一一对应。
\end{solution}

\begin{problem}
    考虑一双原子链的晶格振动，链上最近邻原子间的力常数交错地等于 $c$ 和 $10c$。
令两种原子的质量相等，并且最近邻的间距是 $a/2$，
试求在 $k = 0$ 和 $k = \pi/a$ 处的 $\omega(k)$，
并粗略画出色散关系。本题模拟双原子分子晶体，如 $H_2$。
\end{problem}

\begin{solution}
    双原子链色散关系：
\[
\omega^2 = \frac{c_1 + c_2}{m}
\pm \frac{1}{m}\sqrt{c_1^2 + c_2^2 + 2c_1c_2\cos ka}
\]

代入 $c_1 = c,\; c_2 = 10c$：

\[
\omega^2 = \frac{11c}{m}
\pm \frac{c}{m}\sqrt{101 + 20\cos ka}
\]

(1) $k=0$：
\[
\omega = 0,\quad \omega = \sqrt{\frac{22c}{m}}
\]

(2) $k=\pi/a$：
\[
\omega^2 = \frac{11c}{m} \pm \frac{9c}{m}
\Rightarrow
\omega = \sqrt{\frac{2c}{m}},\quad \sqrt{\frac{20c}{m}}
\]

\[
\boxed{
\text{出现声学支和光学支，并存在带隙}
}
\]
\end{solution}

\begin{problem}
    考虑一个全同原子组成的平面方格子，用 $u_{l,m}$ 记第 $l$ 行、第 $m$ 列
的原子垂直于格平面的位移，每个原子质量为 $M$，最近邻原子的力常数为 $c$。

\begin{enumerate}
\item[(a)] 证明运动方程为：
\[
M \frac{d^2 u_{l,m}}{dt^2}
= c \Big[ (u_{l+1,m} + u_{l-1,m} - 2u_{l,m})
+ (u_{l,m+1} + u_{l,m-1} - 2u_{l,m}) \Big]
\]

\item[(b)] 设解的形式为
\[
u_{l,m} = u(0)\exp\left[i(lk_x a + mk_y a - \omega t)\right]
\]
这里 $a$ 是最近邻原子的间距，证明运动方程是可以满足的，如果
\[
\omega^2 M = 2c(2 - \cos k_x a - \cos k_y a)
\]
这就是问题的色散关系。

\item[(c)] 证明独立解存在的 $k$ 空间区域是一个边长为 $\frac{2\pi}{a}$ 的正方形，
这就是平方格子的第一布里渊区。构出 $k = k_x$ 而 $k_y = 0$ 时，
和 $k_x = k_y$ 时的 $\omega - k$ 图。

\item[(d)] 对于 $ka \ll 1$，证明
\[
\omega = \left(\frac{ca^2}{M}\right)^{1/2} (k_x^2 + k_y^2)^{1/2}
= \left(\frac{ca^2}{M}\right)^{1/2} k
\]
\end{enumerate}
\end{problem}

\begin{solution}
    (a) 最近邻作用：
\[
M \ddot{u}_{l,m}
= c[(u_{l+1,m} + u_{l-1,m} - 2u_{l,m})
+ (u_{l,m+1} + u_{l,m-1} - 2u_{l,m})]
\]

(b) 代入平面波：
\[
u_{l,m} = u_0 e^{i(lk_x a + mk_y a - \omega t)}
\]

得到：
\[
\omega^2 M = 2c(2 - \cos k_x a - \cos k_y a)
\]

(c) 第一布里渊区：
\[
\boxed{k_x, k_y \in \left[-\frac{\pi}{a}, \frac{\pi}{a}\right]}
\]

(d) 长波极限：
\[
\cos x \approx 1 - \frac{x^2}{2}
\]

\[
\omega^2 = \frac{ca^2}{M}(k_x^2 + k_y^2)
\Rightarrow
\boxed{\omega = \left(\frac{ca^2}{M}\right)^{1/2} k}
\]
\end{solution}
\end{document}